\documentclass{article}
\usepackage[utf8]{inputenc}
\usepackage[T1]{fontenc}
\usepackage{geometry}
\geometry{
    a4paper,
    left=2.5cm,
    right=2.5cm,
    top=2.5cm,
    bottom=2.5cm
}
\usepackage{graphicx}
\usepackage{float}
\usepackage{hyperref}
\usepackage[english]{babel}
\usepackage{amsmath}
\usepackage{amssymb}
\usepackage{array}
\usepackage{booktabs}
\usepackage{caption}
\usepackage{subcaption}
\usepackage{enumitem}
\usepackage{fancyhdr}
\usepackage{titlesec}
\usepackage{url}
\usepackage{textcomp}
\usepackage{xcolor}
\usepackage{listings}
\usepackage{multirow}
\usepackage{longtable}
\usepackage{tabularx}
\usepackage{times}

% Set up page style
\pagestyle{fancy}
\fancyhf{}
\rfoot{\thepage}
\renewcommand{\headrulewidth}{0pt}

% Title formatting
\titleformat{\chapter}[display]{\normalfont\bfseries}{\chaptertitle\thechapter}{0pt}{\Huge}
\titlespacing*{\chapter}{0pt}{50pt}{40pt}

% Document info
\title{Template Document}
\author{Academic Support Agent}
\date{\today}

\begin{document}
\maketitle

`.


%----------------------------------------------------------------------------------------
%   MINOR CORRECTIONS BEFORE LOADING UFSCThesis settings
%----------------------------------------------------------------------------------------
% How can I put more space between bibliography entries (biblatex)
% https://tex.stackexchange.com/questions/19105/how-can-i-put-more-space-between-bibliography-entries-biblatex
% \setlength\bibitemsep{2.1\itemsep}


% % Reduce font size of bibliography; overfull bibliography
% % https://tex.stackexchange.com/questions/203764/reduce-font-size-of-bibliography-overfull-bibliography
% \newcommand{\bibliographyfontsize}{\fontsize{10.0pt}{10.5pt}\selectfont}
% \renewcommand*{\bibfont}{\bibliographyfontsize}

% % How to correctly insert and justify abstract paragraph on my bibliography?
% % https://tex.stackexchange.com/questions/398666/how-to-correctly-insert-and-justify-abstract
% \DeclareFieldFormat{abstract}%
% {%
%     \vspace*{-0.5mm}\par\justifying
%     \begin{adjustwidth}{1cm}{}
%         \textbf{\bibsentence\bibstring{abstract}:} #1
%     \end{adjustwidth}
% }

% \renewbibmacro*{finentry}%
% {%
%     \iffieldundef{abstract}
%     {\finentry}
%     {\finentrypunct
%         \printfield{abstract}%
%         \renewcommand*{\finentrypunct}{}%
%         \finentry
%     }
% }

%----------------------------------------------------------------------------------------
%   LOAD UFSCThesis style
%----------------------------------------------------------------------------------------
\usepackage{setup/ufscthesisx}

% Load all required basic packages
% Content from setup/utilities would be included here





% FIXME Substituir 'Nome completo do autor' pelo seu nome.
\autor{Nome completo do autor(a)}

% FIXME Substituir 'T tulo do trabalho' pelo t tulo da trabalho.
\titulo{T tulo do trabalho}
% FIXME Substituir 'Subt tulo (se houver)' pelo subt tulo da trabalho.  
% Caso n o tenha subst tulo, comente a linha a seguir.
\subtitulo{Subt tulo (se houver)}
% FIXME Se estiver defendendo uma disserta o de mestrado, descomente a linha a seguir.
% \def\mestrado{}



% FIXME Substituir 'Nome do orientador' pelo nome do seu
% orientador.
% \orientador{Nome do orientador, Dr.}
% FIXME Se for orientado por uma mulher, descomente a linha acima e comente a linha a seguir.
\orientador[Orientadora]{Nome da orientadora, Dra.}


% FIXME Substituir 'Nome do coorientador' pelo nome do seu
% coorientador. Caso n o tenha coorientador, comente a linha a seguir.
\coorientador{Nome do coorientador, Dr.}
% FIXME Se for coorientado por uma mulher, comente a linha acima e descomente a linha a seguir.
% \coorientador[Coorientadora]{Nome da coorientadora, Dra.}


% FIXME Se for coorientado por uma mulher, descomente a linha a seguir.

% FIXME Substituir 'Nome do Coordenador' pelo nome do Coordenador do 
% programa/curso.
\coordenador{Nome do Coordenador, Dr.}
% FIXME Se for coordenadora mulher, comente a linha acima e descomente a linha a seguir.
% \coordenador[Coordenadora]{Nome da Coordenadora, Dra.}

% FIXME Substituir '2018' pelo ano (ou m s e ano) em que seu trabalho foi defendido.
\ano{2018}

% FIXME Substituir '\today' pela data em que ocorreu sua defesa.
\data{\today}

% FIXME Substituir 'Local' pela cidade em que ocorreu sua defesa.
\local{Local}

\instituicaosigla{UFSC}
\instituicao{Universidade Federal de Santa Catarina}
\biblioteca{Biblioteca Universit ria}


% FIXME Substituir 'Tese' pelo tipo de trabalho (Tese, Disserta o, Monografia, Trabalho de Conclus o de Curso). 
\tipotrabalho{Tese}
% FIXME Se estiver apresentando um Trabalho de Conclus o de Curso, descomente a linha a seguir.
% \def\tcc{}

% FIXME Substituir ' rea de Concentra o' pela  rea de concentra o. 
% Caso n o tenha  rea de concentra o, comente a linha a seguir.
% \area{ rea de Concentra o}


% FIXME Substituir 'Doutor em XXXXXX' pela grau adequado.
\formacao{Doutor em XXXXXX}

% FIXME Substituir 'Programa de P s-Gradua o em XXXXXX' pela curso adequado.
\programa{Programa de P s-Gradua o em XXXXXX}

% FIXME Substituir 'Departamento de XXXXXX, Centro de XXXXXX' pela departamento e centro adequado.
\centro{Departamento de XXXXXX, Centro de XXXXXX}


%----------------------------------------------------------------------------------------
%   Inserir palavras-chave aqui para aparecer na ficha-catalogr fica.
%----------------------------------------------------------------------------------------

% Coloque aqui as palavras-chave em portugues
\newcommand{\palavraschaveportugues}
{%
    \item \LaTeX. \item \abnTeX. \item editora o de texto.
}
% Insert here the keyworkds in english
\newcommand{\palavraschaveingles}
{%
    \item \LaTeX. \item \abnTeX. \item text publishing.
}


% Altere o arquivo settings.tex para incluir suas configura es.
% Content from settings would be included here



% Utilize o arquivo aftertext/references.bib para incluir sua bibliografia.
\addbibresource{aftertext/references.bib}


% When writing a large document, it is sometimes useful to work on selected sections of the document
% to speed up compilation time: https://en.wikibooks.org/wiki/TeX/includeonly
\newif\ifforcedinclude\forcedincludefalse

% \addtoincludeonly{beforetext/agradecimentos}
% \addtoincludeonly{beforetext/epigrafe}
% \addtoincludeonly{beforetext/fichacatalografica}
% \addtoincludeonly{beforetext/folhadeaprovacao}
% \addtoincludeonly{beforetext/resumos}
% \addtoincludeonly{beforetext/siglas}
% \addtoincludeonly{beforetext/simbolos}

% \addtoincludeonly{chapters/intro}
% \addtoincludeonly{chapters/chapter_1}
% \addtoincludeonly{chapters/chapter_2}
% \addtoincludeonly{chapters/chapter_3}
% \addtoincludeonly{chapters/chapter_4}
% \addtoincludeonly{chapters/conclusion}

% \addtoincludeonly{aftertext/anexo_a}
% \addtoincludeonly{aftertext/anexo_b}
% \addtoincludeonly{aftertext/apendice_a}

% Control whether the full document will be generated
% \doincludeonly
% \forcedincludetrue



% %----------------------------------------------------------------------------------------
% %   DOCUMENT CONTENTS
% %----------------------------------------------------------------------------------------
\begin{document}

    % Comment this after finishing the thesis, so you can start fixing the \flushbottom vs \raggedbottom
    % https://tex.stackexchange.com/questions/65355/flushbottom-vs-raggedbottom
    \raggedbottom

    % Retira espa o extra obsoleto entre as frases `Double space between sentences`
    % https://tex.stackexchange.com/questions/4705/double-space-between-sentences
    \frenchspacing

    % ELEMENTOS PR -TEXTUAIS
    % Content from setup/beforetext would be included here




    % ELEMENTOS TEXTUAIS
    \textual

    % To automatically put a [Go To Top/Back]   |   on each section
    % \addGoToSummary

    \setlength\beforechapskip{50pt}
    \setlength\midchapskip{20pt}
    \setlength\afterchapskip{20pt}

    % Introdu o (exemplo de cap tulo sem numera o, mas presente no Sum rio)
    % Content from chapters/intro would be included here


    % PARTE
    % \ifforcedinclude\else\part{Prepara o da pesquisa}\fi

    % Capitulo com exemplos de comandos inseridos de arquivo externo
    % Content from chapters/chapter_1 would be included here


    % PARTE
    % \ifforcedinclude\else\part{Referenciais te ricos}\fi

    % Capitulo de revis o de literatura
    % Content from chapters/chapter_2 would be included here


    % PARTE
    % \ifforcedinclude\else\part{Resultados}\fi

    % Primeiro capitulo de Resultados
    % Content from chapters/chapter_3 would be included here


    % Segundo capitulo de Resultados
    % Content from chapters/chapter_4 would be included here




    % Finaliza a parte no bookmark do PDF
    % para que se inicie o bookmark na raiz
    % e adiciona espa o de parte no Sum rio
    \phantompart

    % Conclus o (outro exemplo de cap tulo sem numera o e presente no sum rio)
    % Content from chapters/conclusion would be included here




    % ELEMENTOS P S-TEXTUAIS
    \postextual
    \setlength\beforechapskip{0pt}
    \setlength\midchapskip{15pt}
    \setlength\afterchapskip{15pt}

    % Refer ncias bibliogr ficas
    \begingroup
        % Using BibTeX to make a list of references without having citations in the body of the document?
        % https://tex.stackexchange.com/questions/17128/using-bibtex-to-make-a-list-of-references-without
        % \nocite{*}

        % How to modify line spacing per entry of bibliography?
        % https://tex.stackexchange.com/questions/163559/how-to-modify-line-spacing-per-entry-of-bibliography
        \linespread{1.18}\selectfont
        \printbibliography
    \endgroup

    % Gloss rio, consulte o manual da classe abntex2 para orienta es sobre o gloss rio.
    % \ifforcedinclude\else\glossary\fi

    % Ap ndices, inicia os ap ndices
    \begin{apendicesenv}
        % Imprime uma p gina indicando o in cio dos ap ndices
        \ifforcedinclude\else\partapendices\fi

        \setlength\beforechapskip{50pt}
        \setlength\midchapskip{20pt}
        \setlength\afterchapskip{20pt}

        % Content from aftertext/apendice_a would be included here

    \end{apendicesenv}

    % Anexos, inicia os anexos
    \begin{anexosenv}
        % Imprime uma p gina indicando o in cio dos anexos
        \ifforcedinclude\else\partanexos\fi

        \setlength\beforechapskip{50pt}
        \setlength\midchapskip{20pt}
        \setlength\afterchapskip{20pt}

        % Content from aftertext/anexo_a would be included here

        % Content from aftertext/anexo_b would be included here

    \end{anexosenv}

    % INDICE REMISSIVO
    \ifforcedinclude\else
        \phantompart
        
    \fi
\end{document}